\subsection{Forty-six ordinary genus-two moduli classes over \texorpdfstring{$\F_4$}{F4}}

\begin{claim}{A.C2.22}{Proved}
Among the $48$ geometric ordinary genus-two moduli
classes fixed by fourth-power Frobenius, the $46$
classes in the table below satisfy A1.
Equivalently, every ordinary genus-two curve over
$\F_4$ is covered by this conclusion unless its
normalized unordered coefficient triple has one
of the two symmetric polynomials
\[
 Z^3+\alpha Z^2+\alpha^2,\qquad
 Z^3+\alpha^2Z^2+\alpha,\qquad \alpha^2+\alpha+1=0.              \tag{C2.65}
\]
The assertion includes every finite extension,
every finitely generated form of the $46$
positive geometric classes, and every hyperbolic
open subcurve. The two classes in (C2.65) are not
decided by this enumeration; A.C2.25 below settles
both by a different complete primitive theta packet.
\end{claim}
\begin{proof}
\proofparagraph{The exact list of geometric classes.}
Over an algebraic closure every ordinary genus-two
curve in characteristic two has a reduced form
\[
 C(a,b,c):\quad y^2+y=\phi(x)=ax+\frac bx+\frac c{x+1},
                        \qquad abc\ne0.                        \tag{C2.66}
\]
Indeed the Artin--Schreier two-rank formula gives
three branch points, and the genus-two different
formula forces all three reduced pole orders
to be one. Normalize their coordinates to
$0,1,\infty$; a constant is removed after a finite
extension. These are the same normalizations
used in A.C2.1--A.C2.18.

Any geometric isomorphism preserves the unique
hyperelliptic pencil and permutes those three
branch coordinates. The changes $x\mapsto x+1$
and $x\mapsto1/x$ generate the permutations
of $a,b,c$, up to a constant.
Over the identity base map a generator change
has the form $y\mapsto e(x)y+v(x)$.
The coefficient of $y$ requires $e^2+e=0$,
so $e=1$. A nonconstant $v$ has a pole, and
$v^2+v$ then has even positive pole order,
which cannot equal a rational function with
only simple poles. Thus $v$ is constant.
Consequently the unordered triple $\{a,b,c\}$
is exactly the geometric moduli datum.

For a curve over $\F_4$, fourth-power Frobenius
preserves this triple. The possible triples
are precisely the root multisets of monic
cubics over $\F_4$ with nonzero constant term.
Their number is $4\cdot4\cdot3=48$.
There are ten split triples, eighteen triples
with splitting type $1+2$, and twenty irreducible
cubic triples. Their coefficients lie respectively
in $\F_4,\F_{16},\F_{64}$.
Repeated coefficient roots are allowed; the
three branch coordinates remain distinct.
Only containment of every actual $\F_4$ curve
in this list is needed, so no separate descent
existence theorem for stable moduli points is
being assumed.

Encode $\F_4$ by $0,1,2=\alpha,3=\alpha^2$.
A key $s_1s_2s_3$ means $Z^3+s_1Z^2+s_2Z+s_3$.
In the following table $q$ is the coefficient-field
size, $n$ is a theta-class annihilator, and $e$
is the degree of the chosen coordinate factor
over that field. The omitted keys are exactly
$203,302$.
\begin{center}
\begin{tabular}{cccc@{\qquad}cccc}
\toprule
Key & $q$ & $n$ & $e$ & Key & $q$ & $n$ & $e$\\
\midrule
001 & 4 & 7 & 1 & 133 & 4 & 15 & 2\\
002 & 64 & 27 & 1 & 201 & 64 & 3 & 1\\
003 & 64 & 27 & 1 & 202 & 16 & 23 & 1\\
011 & 64 & 5 & 1 & 211 & 16 & 3 & 1\\
012 & 16 & 7 & 1 & 212 & 4 & 9 & 2\\
013 & 16 & 7 & 1 & 213 & 64 & 5 & 1\\
021 & 64 & 13 & 1 & 221 & 16 & 3 & 1\\
022 & 16 & 15 & 1 & 222 & 64 & 9 & 1\\
023 & 16 & 5 & 1 & 223 & 4 & 7 & 2\\
031 & 64 & 13 & 1 & 231 & 4 & 3 & 1\\
032 & 16 & 5 & 1 & 232 & 64 & 19 & 1\\
033 & 16 & 15 & 1 & 233 & 64 & 11 & 1\\
101 & 64 & 3 & 1 & 301 & 64 & 3 & 1\\
102 & 16 & 49 & 3 & 303 & 16 & 23 & 1\\
103 & 16 & 49 & 3 & 311 & 16 & 3 & 1\\
111 & 4 & 3 & 1 & 312 & 64 & 5 & 1\\
112 & 64 & 7 & 1 & 313 & 4 & 9 & 2\\
113 & 64 & 7 & 1 & 321 & 4 & 3 & 1\\
121 & 16 & 3 & 1 & 322 & 64 & 11 & 1\\
122 & 4 & 15 & 2 & 323 & 64 & 19 & 1\\
123 & 64 & 11 & 1 & 331 & 16 & 3 & 1\\
131 & 16 & 3 & 1 & 332 & 4 & 7 & 2\\
132 & 64 & 11 & 1 & 333 & 64 & 9 & 1\\
\bottomrule
\end{tabular}
\end{center}
The distinct displayed annihilators are
\[
                  3,5,7,9,11,13,15,19,23,27,49.
\]
For every one of them, direct multiplication
shows $\langle2,-1\rangle=(\Z/n)^\times$.
This finite table is a list of positive
certificates, not a conclusion from negative
searches at other orders.

\proofparagraph{The fixed exact witnesses.}
The complete data in
\path{code/A_theta_F4_moduli_witnesses.json}
give, for every row, the exact coefficient
triple, an irreducible coordinate polynomial
$q_0(T)$ over its coefficient field, a root
$t\ne0,1$, and $\delta=\phi(t)$.
They also give polynomials $A(x),B(x)$ with
\[
 \deg A\le d,\quad B\text{ monic of degree }d=(n-3)/2,
 \quad A(t)=0,\quad B(t)B(0)B(1)\ne0.
\]
With $u=x(x+1)$ and
$F=b+(b+c)x+ax^2+ax^3=u\phi(x)$, each record
verifies the full polynomial identity
\[
            u(A^2+AB)+B^2(F+\delta u)=a(x+t)^n.                \tag{C2.67}
\]
The coefficient fields have respective binary
moduli $7,19,67$, namely
\[
 w^2+w+1,\qquad w^4+w+1,\qquad w^6+w+1.
\]
Elements of the coordinate field are encoded
by consecutive blocks of $2,4,$ or $6$ bits
for their coefficients in $1,T,\ldots,T^{e-1}$.
All chosen coordinate factors have degree
at most three. The largest coordinate fields
are $\F_{4096}$ in the two degree-$49$ rows.
The auxiliary $\beta$ below may require a
further quadratic constant extension.

The exact verifier
\path{code/A_theta_F4_moduli_witnesses.py}
uses ordinary polynomial coefficient lists
and imports none of the determinant engine.
It checks the finite-field operations and
irreducibility, solves the original local
jet equations, verifies all those jets, and
multiplies (C2.67) in full. Every one of the
$46$ fixed records passes. Thus the positive
certificates can be checked without any
large determinant or inference from numerical
point counts.

Choose $\beta^2+\beta=\delta$ and put
$R=(t,\beta)$, $h=A+B(y+\beta)$.
Its norm is $a(x+t)^n/u$. At $R$, $h=0$
and its conjugate has nonzero value $B(t)$.
The $x$-map is \'etale there, and the norm
has its only zero at $t$, of order $n$.
The values $B(0),B(1)\ne0$ give simple poles
at $W_0,W_1$. At infinity, $x^dy$ has the
unique largest pole order $2d+1=n-2$.
Therefore
\[
 \begin{aligned}
 \operatorname{div}(h)&=nR-W_0-W_1-(n-2)W_\infty,\\
 \operatorname{div}(h/u^{(n-1)/2})&=n(R-B_\theta).
 \end{aligned}                                                 \tag{C2.68}
\]
The theta class $L=\mathcal O(R-B_\theta)$
is nonzero: otherwise $2R$ is canonical,
forcing $R$ to be Weierstrass, contrary to
$t\ne0,1,\infty$.
Its actual order $n_0>1$ divides $n$.
The reduction map $(\Z/n)^\times\to(\Z/n_0)^\times$
is surjective, so the generation property
of $2,-1$ descends to $n_0$.
The compact criterion in A.C2.17, or A.C2.19,
proves A1. In particular, exactness of the
listed annihilator is unnecessary, and no
negative theta scan enters this implication.
The cover, isomorphism and parameters descend
after finite extension of the field of
definition; A.3 and the geometric contradiction
then give all stated form and extension
consequences. Equality on an open subcurve
maps to equality on its compactification.

\proofparagraph{The largest two witnesses are small finite-field identities.}
For key $102$, take $\F_{16}$ with
$w^4+w+1=0$, $w$ encoded by $2$, and
$(a,b,c)=(6,8,15)$. Its coordinate polynomial is
\[
                         q_0(T)=T^3+2T^2+4T+14.
\]
For key $103$, take $(a,b,c)=(7,10,12)$ and
\[
                         q_0(T)=T^3+5T^2+2T+13.
\]
Neither cubic has a root in $\F_{16}$, so both
are irreducible. The following matrices list
the coefficients of $A$ and $B$ in increasing
powers of $x$, read row by row; integers use
four-bit coefficient blocks in the indicated
$\F_{4096}$.
For key $102$, $\delta=1663$, and
\[
\begin{aligned}
\operatorname{coeff}(A)&=\left[\begin{matrix}
2967 & 3094 & 2226 & 1784 & 3122 & 1282\\
2330 & 2961 & 160 & 1454 & 3913 & 1774\\
2664 & 2041 & 396 & 120 & 2115 & 3265\\
3418 & 2933 & 2951 & 2846 & 2919 & 1547
\end{matrix}\right],\\[3pt]
\operatorname{coeff}(B)&=\left[\begin{matrix}
1357 & 1974 & 2300 & 2749 & 1451 & 1952\\
3656 & 3667 & 1719 & 3593 & 4043 & 3650\\
2204 & 3665 & 3643 & 1821 & 57 & 3133\\
1148 & 362 & 3621 & 1909 & 111 & 1
\end{matrix}\right].
\end{aligned}
\]
For key $103$, $\delta=1900$, and
\[
\begin{aligned}
\operatorname{coeff}(A)&=\left[\begin{matrix}
3766 & 2071 & 2789 & 1994 & 2117 & 773\\
2847 & 3761 & 240 & 1021 & 3115 & 2013\\
3962 & 1739 & 424 & 106 & 2596 & 2177\\
2367 & 3683 & 3750 & 3613 & 3702 & 1806
\end{matrix}\right],\\[3pt]
\operatorname{coeff}(B)&=\left[\begin{matrix}
809 & 1767 & 2760 & 4073 & 1022 & 1776\\
3370 & 3380 & 2022 & 3339 & 3214 & 3365\\
2744 & 3377 & 3406 & 1561 & 75 & 2121\\
616 & 383 & 3411 & 1635 & 124 & 1
\end{matrix}\right].
\end{aligned}
\]
These two identities alone verify the largest
positive entries without recomputing their
$25$ maximal minors. The separate discovery
scan also excluded theta seven-torsion there,
so the displayed class actually has exact
order $49$; this stronger fact is not needed
for the table's point-section conclusion.

\proofparagraph{The actual remaining pair.}
For key $203$, the normalized coefficients are
\[
 \F_{16}=\F_2[w]/(w^4+w+1),\quad \alpha=w^2+w,
 \qquad(a,b,c)=(1,w^3+1,w^3+w^2+w)=(1,9,14).
\]
Fourth-power Frobenius swaps $b,c$, inducing
branch permutation $0\leftrightarrow1$.
An explicit $\F_4$ model is
\[
              v^2+v=z+\frac{\alpha^2(z+1)}{z^2+z+\alpha}.
                                                               \tag{C2.69}
\]
Indeed put $x=z+w$ and $y=v+k$ with
$k^2+k=w$, which has a solution in $\F_{16}$.
The rational function transforms exactly
as displayed. Key $302$ is the coefficientwise
square conjugate, with $\alpha$ replaced by
$\alpha^2$ and normalized coefficients
$(1,11,13)$.
These two classes have not been decided by
the table. Bounded further scans and partial
primitive-orbit hits are kept separate from
this positive theorem; they yield neither
an A1 counterexample nor a universal bound.

The exact geometric discovery criterion is
proved under A.C2.18 and in
\path{audit/A_theta_Hankel_scan.md}.
Original separate-determinant records are
preserved in
\path{code/A_theta_F4_moduli_original_minors.json}.
An augmented-Bareiss adjugate implementation
was compared with $189$ identical exact gcds
from those records; the check is
\path{code/A_theta_F4_moduli_adjugate_check.json}.
The fixed norm certificates are an independent
positive verification and are sufficient for
the present theorem.
A wholly independent SymPy quotient-ring implementation,
\path{code/A_F4_norm_root_independent.py}, also checks
all $46$ identities, field and coordinate irreducibility,
local values, moduli keys and unit-group generation;
its adjacent JSON records PASS. The complete proof and
scope record are
\path{audit/A_F4_moduli_46_candidate.md}.
Universal A1 remains A.30.
\end{proof}
